% ============================================================
% Question Bank: ch01-07 Proportional Reasoning with Scale Models
% Topic Title: Proportional Reasoning with Scale Models
% CCSS: Supplementary
% Questions: 30 (18 MC + 7 SA + 5 Graphical)
% ============================================================

% --- Multiple Choice Questions (18) ---

\begin{practiceQuestion}{1-7-q01}{mc}
\begin{questionText}
A model car is $10$ cm long. The scale is $1 : 40$. What is the actual length of the car?
\end{questionText}
\choiceA{$40$ cm}
\choiceB{$400$ cm}
\choiceC{$4$ cm}
\choiceD{$4{,}000$ cm}
\correctAnswer{B}
\explanation{Multiply model length by scale factor: $10 \times 40 = 400$ cm.}
\end{practiceQuestion}

\begin{practiceQuestion}{1-7-q02}{mc}
\begin{questionText}
A building is $60$ m tall. A model uses the scale $1 : 200$. How tall is the model?
\end{questionText}
\choiceA{$30$ cm}
\choiceB{$3$ cm}
\choiceC{$300$ cm}
\choiceD{$0.3$ cm}
\correctAnswer{A}
\explanation{Divide actual by scale factor: $60 \div 200 = 0.3$ m $= 30$ cm.}
\end{practiceQuestion}

\begin{practiceQuestion}{1-7-q03}{mc}
\begin{questionText}
A model airplane wing is $20$ cm long. The real wing is $16$ m long. What is the scale factor?
\end{questionText}
\choiceA{$1 : 8$}
\choiceB{$1 : 80$}
\choiceC{$1 : 800$}
\choiceD{$1 : 20$}
\correctAnswer{B}
\explanation{Convert: $16$ m $= 1{,}600$ cm. Scale factor: $1{,}600 \div 20 = 80$. Scale is $1 : 80$.}
\end{practiceQuestion}

\begin{practiceQuestion}{1-7-q04}{mc}
\begin{questionText}
On a scale model, $1$ cm represents $25$ m. If a bridge on the model is $6$ cm long, what is the actual length of the bridge?
\end{questionText}
\choiceA{$125$ m}
\choiceB{$150$ m}
\choiceC{$31$ m}
\choiceD{$600$ m}
\correctAnswer{B}
\explanation{Multiply: $6 \times 25 = 150$ m.}
\end{practiceQuestion}

\begin{practiceQuestion}{1-7-q05}{mc}
\begin{questionText}
A model train uses a $1 : 50$ scale. If the real train is $15$ m long, how long is the model?
\end{questionText}
\choiceA{$30$ cm}
\choiceB{$3$ cm}
\choiceC{$75$ cm}
\choiceD{$300$ cm}
\correctAnswer{A}
\explanation{Divide: $15 \div 50 = 0.3$ m $= 30$ cm.}
\end{practiceQuestion}

\begin{practiceQuestion}{1-7-q06}{mc}
\begin{questionText}
A map uses a scale of $1 : 10{,}000$. Two cities are $5$ cm apart on the map. What is the actual distance?
\end{questionText}
\choiceA{$50$ m}
\choiceB{$500$ m}
\choiceC{$5{,}000$ m}
\choiceD{$50{,}000$ m}
\correctAnswer{B}
\explanation{Multiply: $5 \times 10{,}000 = 50{,}000$ cm $= 500$ m.}
\end{practiceQuestion}

\begin{practiceQuestion}{1-7-q07}{mc}
\begin{questionText}
A room is $8$ m by $5$ m. A scale model uses $1 : 100$. What are the model's dimensions?
\end{questionText}
\choiceA{$80$ cm by $50$ cm}
\choiceB{$8$ cm by $5$ cm}
\choiceC{$0.8$ cm by $0.5$ cm}
\choiceD{$4$ cm by $2.5$ cm}
\correctAnswer{B}
\explanation{Divide each dimension by $100$: $800 \div 100 = 8$ cm and $500 \div 100 = 5$ cm.}
\end{practiceQuestion}

\begin{practiceQuestion}{1-7-q08}{mc}
\begin{questionText}
A $1 : 250$ model of a skyscraper is $48$ cm tall. What is the actual height of the building?
\end{questionText}
\choiceA{$12$ m}
\choiceB{$120$ m}
\choiceC{$1{,}200$ m}
\choiceD{$12{,}000$ m}
\correctAnswer{B}
\explanation{Multiply: $48 \times 250 = 12{,}000$ cm $= 120$ m.}
\end{practiceQuestion}

\begin{practiceQuestion}{1-7-q09}{mc}
\begin{questionText}
A statue is $3$ m tall. An artist makes a model that is $15$ cm tall. What scale did the artist use?
\end{questionText}
\choiceA{$1 : 5$}
\choiceB{$1 : 15$}
\choiceC{$1 : 20$}
\choiceD{$1 : 200$}
\correctAnswer{C}
\explanation{Convert: $3$ m $= 300$ cm. Scale: $300 \div 15 = 20$. The scale is $1 : 20$.}
\end{practiceQuestion}

\begin{practiceQuestion}{1-7-q10}{mc}
\begin{questionText}
On a $1 : 500$ scale model, a parking lot is $4$ cm wide. How wide is the real parking lot?
\end{questionText}
\choiceA{$200$ m}
\choiceB{$20$ m}
\choiceC{$2$ m}
\choiceD{$2{,}000$ m}
\correctAnswer{B}
\explanation{Multiply: $4 \times 500 = 2{,}000$ cm $= 20$ m.}
\end{practiceQuestion}

\begin{practiceQuestion}{1-7-q11}{mc}
\begin{questionText}
A model ship is $35$ cm long. The scale is $1 : 300$. What is the actual length of the ship?
\end{questionText}
\choiceA{$10.5$ m}
\choiceB{$105$ m}
\choiceC{$1{,}050$ m}
\choiceD{$35$ m}
\correctAnswer{B}
\explanation{Multiply: $35 \times 300 = 10{,}500$ cm $= 105$ m.}
\end{practiceQuestion}

\begin{practiceQuestion}{1-7-q12}{mc}
\begin{questionText}
A swimming pool is $25$ m long. A scale model is $50$ cm long. What is the scale?
\end{questionText}
\choiceA{$1 : 25$}
\choiceB{$1 : 50$}
\choiceC{$1 : 500$}
\choiceD{$1 : 5$}
\correctAnswer{B}
\explanation{Convert: $25$ m $= 2{,}500$ cm. Scale: $2{,}500 \div 50 = 50$. The scale is $1 : 50$.}
\end{practiceQuestion}

\begin{practiceQuestion}{1-7-q13}{mc}
\begin{questionText}
A $1 : 75$ model truck is $8$ cm long and $3$ cm wide. What is the actual width of the truck?
\end{questionText}
\choiceA{$2.25$ m}
\choiceB{$6$ m}
\choiceC{$225$ m}
\choiceD{$0.225$ m}
\correctAnswer{A}
\explanation{Width: $3 \times 75 = 225$ cm $= 2.25$ m.}
\end{practiceQuestion}

\begin{practiceQuestion}{1-7-q14}{mc}
\begin{questionText}
An architect's model of a house is built at scale $1 : 50$. The real house is $12$ m long and $9$ m wide. What is the length of the model?
\end{questionText}
\choiceA{$2.4$ cm}
\choiceB{$24$ cm}
\choiceC{$18$ cm}
\choiceD{$6$ cm}
\correctAnswer{B}
\explanation{Convert: $12$ m $= 1{,}200$ cm. Model length: $1{,}200 \div 50 = 24$ cm.}
\end{practiceQuestion}

\begin{practiceQuestion}{1-7-q15}{mc}
\begin{questionText}
Two students build models of the same tower. Student A uses scale $1 : 100$; Student B uses scale $1 : 200$. The tower is $80$ m tall. How much taller is Student A's model than Student B's?
\end{questionText}
\choiceA{$20$ cm}
\choiceB{$40$ cm}
\choiceC{$80$ cm}
\choiceD{$60$ cm}
\correctAnswer{B}
\explanation{A: $8{,}000 \div 100 = 80$ cm. B: $8{,}000 \div 200 = 40$ cm. Difference: $80 - 40 = 40$ cm.}
\end{practiceQuestion}

\begin{practiceQuestion}{1-7-q16}{mc}
\begin{questionText}
A model helicopter is $18$ cm long. The actual helicopter is $9$ m long. What does $1$ cm on the model represent in real life?
\end{questionText}
\choiceA{$50$ cm}
\choiceB{$0.5$ m}
\choiceC{$5$ m}
\choiceD{$2$ m}
\correctAnswer{B}
\explanation{Scale: $900 \div 18 = 50$. So $1$ cm represents $50$ cm $= 0.5$ m.}
\end{practiceQuestion}

\begin{practiceQuestion}{1-7-q17}{mc}
\begin{questionText}
A football field is $100$ m long. On a $1 : 500$ scale model, how long is the field?
\end{questionText}
\choiceA{$20$ cm}
\choiceB{$5$ cm}
\choiceC{$50$ cm}
\choiceD{$2$ cm}
\correctAnswer{A}
\explanation{Divide: $10{,}000 \div 500 = 20$ cm.}
\end{practiceQuestion}

\begin{practiceQuestion}{1-7-q18}{mc}
\begin{questionText}
A model of the Eiffel Tower is $33$ cm tall. The real tower is $330$ m tall. Which scale was used?
\end{questionText}
\choiceA{$1 : 10$}
\choiceB{$1 : 100$}
\choiceC{$1 : 1{,}000$}
\choiceD{$1 : 10{,}000$}
\correctAnswer{C}
\explanation{Convert: $330$ m $= 33{,}000$ cm. Scale: $33{,}000 \div 33 = 1{,}000$. The scale is $1 : 1{,}000$.}
\end{practiceQuestion}

% --- Short Answer Questions (7) ---

\begin{practiceQuestion}{1-7-q19}{sa}
\begin{questionText}
A model bridge is $14$ cm long. The scale is $1 : 150$. What is the actual length of the bridge in meters?
\end{questionText}
\correctAnswer{$21$ m}
\explanation{Multiply: $14 \times 150 = 2{,}100$ cm $= 21$ m.}
\end{practiceQuestion}

\begin{practiceQuestion}{1-7-q20}{sa}
\begin{questionText}
An airplane is $36$ m long. A model uses scale $1 : 200$. How long is the model in centimeters?
\end{questionText}
\correctAnswer{$18$ cm}
\explanation{Convert: $36$ m $= 3{,}600$ cm. Divide: $3{,}600 \div 200 = 18$ cm.}
\end{practiceQuestion}

\begin{practiceQuestion}{1-7-q21}{sa}
\begin{questionText}
A model castle is $25$ cm wide. The real castle is $50$ m wide. What is the scale of the model?
\end{questionText}
\correctAnswer{$1 : 200$}
\explanation{Convert: $50$ m $= 5{,}000$ cm. Scale: $5{,}000 \div 25 = 200$. The scale is $1 : 200$.}
\end{practiceQuestion}

\begin{practiceQuestion}{1-7-q22}{sa}
\begin{questionText}
A park is $300$ m long and $200$ m wide. A scale model uses $1 : 1{,}000$. What are the dimensions of the model in centimeters?
\end{questionText}
\correctAnswer{$30$ cm by $20$ cm}
\explanation{Length: $30{,}000 \div 1{,}000 = 30$ cm. Width: $20{,}000 \div 1{,}000 = 20$ cm.}
\end{practiceQuestion}

\begin{practiceQuestion}{1-7-q23}{sa}
\begin{questionText}
On a map, $1$ cm represents $5$ km. Two towns are $8$ cm apart on the map. What is the actual distance in kilometers?
\end{questionText}
\correctAnswer{$40$ km}
\explanation{Multiply: $8 \times 5 = 40$ km.}
\end{practiceQuestion}

\begin{practiceQuestion}{1-7-q24}{sa}
\begin{questionText}
A $1 : 60$ model car is $7.5$ cm long. What is the actual length of the car in meters?
\end{questionText}
\correctAnswer{$4.5$ m}
\explanation{Multiply: $7.5 \times 60 = 450$ cm $= 4.5$ m.}
\end{practiceQuestion}

\begin{practiceQuestion}{1-7-q25}{sa}
\begin{questionText}
A model rocket is $40$ cm tall and the real rocket is $120$ m tall. How many centimeters on the model represent $1$ meter on the real rocket?
\end{questionText}
\correctAnswer{$\frac{1}{3}$ cm}
\explanation{Scale: $12{,}000 \div 40 = 300$. So $1$ m on the real rocket is $100 \div 300 = \frac{1}{3}$ cm on the model.}
\end{practiceQuestion}

% --- Graphical Questions (3 GMC + 2 GSA) ---

\begin{practiceQuestion}{1-7-q26}{gmc}
\begin{questionText}
The diagram below shows a model building and its actual building. The model is $6$ cm tall. The scale is $1 : 150$. Which label correctly shows the actual height?

\begin{center}
\begin{tikzpicture}[scale=0.8]
  % Model building
  \draw[thick,funBlue,fill=funBlue!20] (0,0) rectangle (1,1.5);
  \node[below] at (0.5,0) {Model};
  \draw[<->] (-0.3,0) -- (-0.3,1.5) node[midway,left] {$6$ cm};

  % Arrow
  \draw[thick,->] (1.8,0.75) -- (3.2,0.75) node[midway,above] {$\times\;150$};

  % Actual building
  \draw[thick,funRed,fill=funRed!20] (4,0) rectangle (5.5,3.5);
  \node[below] at (4.75,0) {Actual};
  \draw[<->] (5.8,0) -- (5.8,3.5) node[midway,right] {?};
\end{tikzpicture}
\end{center}
\end{questionText}
\choiceA{$90$ m}
\choiceB{$9$ m}
\choiceC{$900$ m}
\choiceD{$0.9$ m}
\correctAnswer{B}
\explanation{Multiply: $6 \times 150 = 900$ cm $= 9$ m.}
\end{practiceQuestion}

\begin{practiceQuestion}{1-7-q27}{gmc}
\begin{questionText}
A scale drawing of a classroom is shown below. Each square on the grid represents $1$ cm. The scale is $1 : 50$. What is the actual length of the classroom?

\begin{center}
\begin{tikzpicture}[scale=0.5]
  \draw[step=1,gray!30,thin] (0,0) grid (12,8);
  \draw[thick,funBlue,fill=funBlue!10] (1,1) rectangle (11,7);
  \draw[<->] (1,0.3) -- (11,0.3) node[midway,below] {$10$ cm};
  \draw[<->] (0.3,1) -- (0.3,7) node[midway,left] {$6$ cm};
  \node at (6,4) {Classroom};
\end{tikzpicture}
\end{center}
\end{questionText}
\choiceA{$50$ m}
\choiceB{$5$ m}
\choiceC{$500$ m}
\choiceD{$10$ m}
\correctAnswer{B}
\explanation{Multiply: $10 \times 50 = 500$ cm $= 5$ m.}
\end{practiceQuestion}

\begin{practiceQuestion}{1-7-q28}{gmc}
\begin{questionText}
The table below shows the model and actual measurements of different objects. Which object has the largest scale factor?

\begin{center}
\begin{tabular}{|l|c|c|}
\hline
\textbf{Object} & \textbf{Model} & \textbf{Actual} \\
\hline
Car & $10$ cm & $5$ m \\
\hline
Boat & $8$ cm & $24$ m \\
\hline
House & $12$ cm & $6$ m \\
\hline
\end{tabular}
\end{center}
\end{questionText}
\choiceA{Car}
\choiceB{Boat}
\choiceC{House}
\choiceD{Car and House (tied)}
\correctAnswer{B}
\explanation{Car: $500 \div 10 = 50$. Boat: $2{,}400 \div 8 = 300$. House: $600 \div 12 = 50$. The boat has the largest scale factor ($1 : 300$).}
\end{practiceQuestion}

\begin{practiceQuestion}{1-7-q29}{gsa}
\begin{questionText}
The scale drawing below shows a garden. Each square represents $1$ cm on the drawing. The scale is $1 : 200$. Find the actual perimeter of the garden in meters.

\begin{center}
\begin{tikzpicture}[scale=0.5]
  \draw[step=1,gray!30,thin] (0,0) grid (10,8);
  \draw[thick,funBlue,fill=green!15] (1,1) rectangle (9,6);
  \draw[<->] (1,0.3) -- (9,0.3) node[midway,below] {$8$ cm};
  \draw[<->] (0.3,1) -- (0.3,6) node[midway,left] {$5$ cm};
  \node at (5,3.5) {Garden};
\end{tikzpicture}
\end{center}
\end{questionText}
\correctAnswer{$52$ m}
\explanation{Model perimeter: $2(8 + 5) = 26$ cm. Actual: $26 \times 200 = 5{,}200$ cm $= 52$ m.}
\end{practiceQuestion}

\begin{practiceQuestion}{1-7-q30}{gsa}
\begin{questionText}
The diagram shows a model and the real object with their measurements. Find the scale of the model.

\begin{center}
\begin{tikzpicture}[scale=0.7]
  % Model
  \draw[thick,funBlue,fill=funBlue!15] (0,0) rectangle (2,1.2);
  \node at (1,0.6) {Model};
  \draw[<->] (0,-0.3) -- (2,-0.3) node[midway,below] {$16$ cm};

  % Arrow
  \node at (3.5,0.6) {$\longrightarrow$};

  % Actual
  \draw[thick,funRed,fill=funRed!15] (5,0) rectangle (8.5,2.5);
  \node at (6.75,1.25) {Actual};
  \draw[<->] (5,-0.3) -- (8.5,-0.3) node[midway,below] {$40$ m};
\end{tikzpicture}
\end{center}
\end{questionText}
\correctAnswer{$1 : 250$}
\explanation{Convert: $40$ m $= 4{,}000$ cm. Divide: $4{,}000 \div 16 = 250$. The scale is $1 : 250$.}
\end{practiceQuestion}
